%%%%%%%%%%%%%%%%%%%%%%%%%%%%%%%%%%%%%%%%%%%%%%%%%%%%%%%%%%%%%%%%%%%%%%
%%                Eigene Makros dieser Masterarbeit                  %%
%%%%%%%%%%%%%%%%%%%%%%%%%%%%%%%%%%%%%%%%%%%%%%%%%%%%%%%%%%%%%%%%%%%%%%

%% Diese Datei sammelt die Makros des Autors. Die Dateien des
%% IAEW-Templates (header.tex, personal_settings.tex) bleiben davon
%% unberuehrt. Eingebunden wird sie in main.tex nach extras/header.tex.
%% TikZ wird bereits in extras/header.tex geladen (\usepackage{pgf, tikz}),
%% ein zusaetzliches \usepackage ist daher nicht noetig.

% Platzhalter fuer noch nicht erstellte Abbildungen.
% Argument: Beschreibung des geplanten Bildinhalts.
% Vor Abgabe muessen alle Vorkommen durch \includegraphics ersetzt sein.
\newcommand{\platzhalterabb}[1]{%
  \begin{tikzpicture}
    \node[draw, dashed, line width=1pt, rounded corners,
          inner sep=1.5em, text width=0.82\textwidth, align=left]
      {\textbf{[PLATZHALTER ABBILDUNG]}\par\vspace{0.6em}\small #1};
  \end{tikzpicture}%
}

% Einheit Euro je Megawatt und Stunde, angelegt am 17.09.2026 auf Wunsch des
% Verfassers (Einheiten als Symbol statt ausgeschrieben). Enger Malpunkt wie
% in den Abbildungen. Im Fliesstext als \EurMWh{} schreiben, sonst
% verschluckt TeX das folgende Leerzeichen.
% BERICHTIGT am 18.09.2026. Die Fassung vom 17.09.2026 lautete
%   \newcommand{\EurMWh}{\si[inter-unit-product=\ensuremath{\mathord{\cdot}}]{\euro/(MW.h)}}
% und verlor den Malpunkt, weil \sisetup{mode=math,detect-all} in
% extras/header.tex Zeile 118 siunitx in den Mathematikmodus zwingt und
% \textperiodcentered dort unzulaessig ist. Das Makro setzte deshalb im
% Fliesstext Euro je Megawattstunde, also die Form, die CLAUDE.md
% Abschnitt 8 ausdruecklich abgrenzt, und erzeugte je Verwendung drei
% Warnungen im Log. Betroffen waren chapter_2.tex, chapter_3.tex,
% chapter_4.tex und attachment_dispatch.tex.
\newcommand{\EurMWh}{\si[inter-unit-product=\text{\textperiodcentered}]{\euro/(MW.h)}}
